\documentclass[a4paper]{article}
\usepackage{amsfonts}
\usepackage{amsmath}
\usepackage{amssymb}
\usepackage{graphicx}     

\begin{document}
  
\noindent \large Auteur: Raj Shah \\
\noindent \large Studierichting: Informatica\\
\noindent \large Oefeningenreeks 3C nr. 12 \\

\medskip

\normalsize

Bewijs: als $0 < b \leq a$, dan geldt $\dfrac{a-b}{a} \leq \ln \dfrac{a}{b} \leq \dfrac{a-b}{b}$\\

\begin{itemize}
	\item $f $ is differentieerbaar in $ ]b, a[$
	\item $f$ is continu in $[b, a]$
\end{itemize}

met $f(x) = \ln x \Rightarrow f'(x) = \dfrac{1}{x}$\\
$\overset{Lagrange}{\Rightarrow} f'(c) = \dfrac{f(a) - f(b)}{a-b} = \dfrac{\ln a - \ln b}{a-b} = \dfrac{\ln \dfrac{a}{b}}{a-b}$\\

$\Rightarrow f'(c) = \dfrac{\ln \dfrac{a}{b}}{a-b} = \dfrac{1}{c} \Leftrightarrow \dfrac{a-b}{c} = \ln \dfrac{a}{b}$ met $\left( b \leq c \leq a \right)$\\

$\Leftrightarrow \dfrac{1}{a} \leq \dfrac{1}{c} \leq \dfrac{1}{b}$\\

$\Leftrightarrow \dfrac{1}{a} \leq \dfrac{\ln \dfrac{a}{b}}{a-b} \leq \dfrac{1}{b}$\\

$\Leftrightarrow  \dfrac{a-b}{a} \leq \ln \dfrac{a}{b} \leq \dfrac{a-b}{b}$\\

\begin{thebibliography}{999}
%\bibitem{a} Referentie
\end{thebibliography}
\end{document} 
